\documentclass[tikz,border=10pt]{standalone}
\usepackage{amsmath,amssymb}
\usepackage{tikz}
\usetikzlibrary{arrows,automata,backgrounds,calendar,chains,decorations,decorations.pathreplacing,decorations.text,fit,matrix,mindmap,patterns,petri,plotmarks,positioning,shadows,shapes,shapes.callouts,shapes.geometric,shapes.misc,spy,trees,calc,intersections,tikzmark,arrows.meta}
\usepackage{xcolor}
\usepackage{pgfplots}
\pgfplotsset{compat=1.16}

% Common math macros from the book
\def\x{\mathbf{x}}
\def\y{\mathbf{y}}
\def\z{\mathbf{z}}
\def\A{\mathbf{A}}
\def\b{\mathbf{b}}
\def\c{\mathbf{c}}
\def\w{\mathbf{w}}
\def\v{\mathbf{v}}
\def\u{\mathbf{u}}
\def\e{\mathbf{e}}
\def\zero{\mathbf{0}}
\def\R{\mathbb{R}}
\def\N{\mathbb{N}}
\def\Z{\mathbb{Z}}
\def\Q{\mathbb{Q}}
\def\st{\text{s.t.}}
\def\vect#1{\mathbf{#1}}
\def\xLP{x^{\mathrm{LP}}}
\def\zLP{z^{\mathrm{LP}}}
\def\xIP{x^{\mathrm{IP}}}

% Common node styles from the book
\tikzset{
    main/.style={circle, draw, fill=gray!20, minimum size=1cm, font=\normalsize},
    dot/.style={circle,fill=blue,minimum size=3pt,inner sep=0pt, outer sep=-1pt},
    vertex/.style={circle,fill=black!25,minimum size=20pt,inner sep=0pt},
    edge/.style={draw,-,thick},
    directed edge/.style={draw,->,>=latex,thick},
}

% Additional macros for 3D/coordinate systems
\newcommand{\xhat}{\hat{\x}}
\newcommand{\yhat}{\hat{\y}}
\newcommand{\zhat}{\hat{\z}}
\newcommand{\ihat}{\hat{\imath}}
\newcommand{\jhat}{\hat{\jmath}}
\newcommand{\khat}{\hat{k}}

\begin{document}
\begin{tikzpicture}
    \draw[->]  (0,0)coordinate(O) -- (8,0) node[anchor=north] {$n$};
    \draw[->]  (0,0) -- (0,5);

\draw[domain=0:8,smooth,variable=\x,blue,name path=c1] plot ({\x},{0.5*\x+2*sin(\x r)+1})node[right]{$c\cdot g(n)$};
\draw[domain=0:8,smooth,variable=\x,blue,name path=c0] plot ({\x},{0.3*0.5*\x+0.3*2*sin(\x r)+0.3*1})node[right]{$g(n)$};
\draw[domain=0:8,smooth,variable=\x,red,name path=c2] plot ({\x},{0.2*\x+0.5*sin(\x r)+2})node[right,black]{$f(n)$};
\fill[red,name intersections={of=c1 and c2}]
    (intersection-1) circle (2pt)
    (intersection-2) circle (2pt)
        (intersection-3) circle (2pt) ;

\draw[dashed] (intersection-3) -- (intersection-3|-O) node[below]{$n_0$};
\end{tikzpicture}
\end{document}
